% ============================================================================
% CAPÍTULO 13 — ORMs e migrações\index{migrações}
% Texto completo (rodada editorial 2)
% ============================================================================
\chapter{ORMs e Migrações: Prisma e Drizzle}
\label{cap:orms}

\begin{caixaconceito}
\textbf{ORMs} (Object-Relational Mappers) traduzem tabelas SQL em objetos
tipados do TypeScript — com segurança de tipos de ponta a ponta. Neste
capítulo você dominará os dois líderes do ecossistema: \textbf{Prisma}
(abstração madura e expressiva) e \textbf{Drizzle} (SQL-first, leve e rápido
em serverless) \cite{prisma-docs,drizzle-docs}.
\end{caixaconceito}

Ao final deste capítulo, você será capaz de:
\begin{objetivos}
  \item Definir schemas de dados com Prisma\index{Prisma} e Drizzle\index{Drizzle};
  \item Criar e aplicar migrações de forma segura;
  \item Consultar e escrever dados com queries tipadas (incluindo relações);
  \item Decidir entre Prisma e Drizzle com critérios técnicos;
  \item Evitar o anti-padrão N+1\index{N+1} e otimizar queries;
  \item Entregar o projeto do capítulo: a API de uma biblioteca.
\end{objetivos}

\section{Por que ORMs}

Escrever SQL manualmente (capítulo~\ref{cap:postgresql}) é poderoso, mas
repetitivo e propenso a erro: um nome de coluna digitado errado vira um erro
\emph{em runtime}, quando o banco reclama — ou pior, uma query que devolve
dados errados.

ORMs geram as queries a partir de \emph{tipos}: o compilador pega erros de
coluna, relação e tipo \emph{antes} de tocar no banco. E as \textbf{migrações}
versionam o schema\index{schema} como o Git versiona o código — o banco evolui com o app, de
forma auditável e reversível.

\begin{caixaconceito}
\textbf{Migração} = arquivo versionado que descreve uma mudança de schema
(``criar tabela usuários'', ``adicionar coluna telefone''). Todos os ambientes
(dev, staging, produção) aplicam a mesma sequência de migrações — nunca mais
``alterar a tabela na mão''. O histórico de migrações é a ``partida de
nascimento'' do banco.
\end{caixaconceito}

\section{Prisma: o padrão de mercado}

O Prisma usa um schema declarativo próprio (\texttt{schema.prisma}) e gera um
\textbf{client tipado} que você importa:

\begin{lstlisting}[style=livro, language=SQL,
  caption={Schema Prisma com relações.}, label={list:prisma-schema}]
generator client {
  provider = "prisma-client-js"
}

datasource db {
  provider = "postgresql"
  url      = env("DATABASE_URL")
}

model Autor {
  id     Int     @id @default(autoincrement())
  nome   String
  livros Livro[]
}

model Livro {
  id      Int    @id @default(autoincrement())
  titulo  String
  ano     Int
  autorId Int
  autor   Autor  @relation(fields: [autorId], references: [id])
  emprestimos Emprestimo[]
}

model Emprestimo {
  id             Int      @id @default(autoincrement())
  livroId        Int
  livro          Livro    @relation(fields: [livroId], references: [id])
  nomeLeitor     String
  dataEmprestimo DateTime @default(now())
  dataDevolucao  DateTime?
}
\end{lstlisting}

\begin{lstlisting}[style=livro, language=Bash,
  caption={Comandos essenciais do Prisma.}, label={list:prisma-comandos}]
npx prisma init
npx prisma migrate dev --name criar-autores-e-livros   # gera e aplica
npx prisma generate                                    # gera o client tipado
npx prisma studio                                      # UI para inspecionar dados
npx prisma migrate deploy                              # aplica migrações pendentes
\end{lstlisting}

\begin{lstlisting}[style=livro, language=TypeScript,
  caption={Queries tipadas com relações.}, label={list:prisma-query}]
import { PrismaClient } from "@prisma/client";

const prisma = new PrismaClient();

// Buscar livros com autor e contagem de empréstimos (sem N+1)
const livros = await prisma.livro.findMany({
  where: { ano: { gte: 2020 } },
  include: {
    autor: true,
    _count: { select: { emprestimos: true } },
  },
  orderBy: { titulo: "asc" },
});

// Criar com relação aninhada
const novo = await prisma.livro.create({
  data: {
    titulo: "Clean Code",
    ano: 2008,
    autor: { connect: { id: 1 } },
  },
});

// Transação: registrar empréstimo + marcar disponibilidade
await prisma.$transaction(async (tx) => {
  await tx.emprestimo.create({ data: { livroId: 2, nomeLeitor: "Ana" } });
  // ... outras operações atômicas
});
\end{lstlisting}

\begin{caixadica}
O \texttt{include} do Prisma gera \emph{um} JOIN eficiente — é a cura oficial
do N+1. E \texttt{prisma \$transaction} garante atomicidade (ACID) entre
operações, como o \texttt{BEGIN}/\texttt{COMMIT} que você viu no SQL.
\end{caixadica}

\section{Drizzle: SQL-first, leve e serverless-friendly}

O Drizzle inverte a filosofia: \emph{SQL com tipos}, sem camada de mágica.
Você escreve queries quase iguais ao SQL — com autocomplete e verificação de
tipos:

\begin{lstlisting}[style=livro, language=TypeScript,
  caption={Schema Drizzle (tabelas tipadas).}, label={list:drizzle-schema}]
import {
  pgTable, serial, text, integer, timestamp, index,
} from "drizzle-orm/pg-core";

export const autores = pgTable("autores", {
  id: serial("id").primaryKey(),
  nome: text("nome").notNull(),
});

export const livros = pgTable(
  "livros",
  {
    id: serial("id").primaryKey(),
    titulo: text("titulo").notNull(),
    ano: integer("ano").notNull(),
    autorId: integer("autor_id")
      .notNull()
      .references(() => autores.id),
  },
  (t) => [index("idx_livros_autor").on(t.autorId)],
);
\end{lstlisting}

\begin{lstlisting}[style=livro, language=TypeScript,
  caption={Query Drizzle próxima do SQL.}, label={list:drizzle-query}]
import { eq, and, gte, sql } from "drizzle-orm";
import { db } from "./db";
import { livros, autores } from "./schema";

// SELECT com JOIN, filtro e agregação
const resultado = await db
  .select({
    titulo: livros.titulo,
    autor: autores.nome,
    emprestimos: sql<number>`count(*)`,
  })
  .from(livros)
  .innerJoin(autores, eq(livros.autorId, autores.id))
  .where(and(gte(livros.ano, 2020)))
  .groupBy(livros.id, autores.nome);

// INSERT, UPDATE, DELETE com tipos
const novo = await db.insert(livros).values({
  titulo: "Entendendo Algoritmos",
  ano: 2017,
  autorId: 1,
}).returning();
\end{lstlisting}

\begin{caixadica}
Drizzle usa \textbf{drizzle-kit} para migrações (\texttt{drizzle-kit migrate})
e \texttt{drizzle-orm} para queries. O ponto forte: bundle \textasciitilde{}50
kB e cold starts rápidos — o que o torna a escolha preferida em serverless
(Vercel, Lambda).
\end{caixadica}

\section{Prisma vs Drizzle: critérios de decisão}

\begin{table}[h]
\centering
\small
\caption{Comparativo Prisma vs Drizzle.}
\label{tab:prisma-drizzle}
\begin{tabularx}{\textwidth}{@{}l X X@{}}
\toprule
\textbf{Critério} & \textbf{Prisma} & \textbf{Drizzle} \\
\midrule
Filosofia & Abstração e DX & SQL-first, camada fina \\
Curva de aprendizado & Suave (schema próprio) & Exige conhecer SQL \\
Bundle/cold start & Maior & \textasciitilde{}50 kB, ideal serverless \\
Queries complexas & Boas, com \texttt{raw} quando preciso & Nativas \\
Migrações & \texttt{prisma migrate} & \texttt{drizzle-kit} \\
Ecossistema & Muito maduro & Crescendo rápido \\
\bottomrule
\end{tabularx}
\end{table}

\begin{caixadica}
Recomendação de mercado: \textbf{Prisma} para a maioria dos produtos (DX e
velocidade de desenvolvimento); \textbf{Drizzle} para serverless de alta
performance e times que querem controle total do SQL. Saber os dois = decisão
técnica de senior.
\end{caixadica}

\section{O anti-padrão N+1}

O erro de performance nº 1 com ORMs: buscar a lista e depois buscar a relação
de \emph{cada} item individualmente. Se você tem 100 livros e acessa
\texttt{livro.autor} de cada um, são \textbf{101 queries} ao banco — cada uma
com latência de rede. O Prisma resolve com \texttt{include}; o Drizzle com
\texttt{JOIN}. \textbf{Sempre confira as queries geradas} (logging do Prisma
ou do Drizzle) — o N+1 é silencioso e só aparece sob carga.

\begin{lstlisting}[style=livro, language=TypeScript,
  caption={Ligando o log de queries (Prisma).}, label={list:orm-log}]
const prisma = new PrismaClient({
  log: ["query", "warn", "error"],
});
// Agora cada query aparece no console — o N+1 fica visível.
\end{lstlisting}

% ============================================================================
\section{Projeto do capítulo: API de uma biblioteca}
\label{sec:projeto13}

\begin{caixaprojeto}
\textbf{Objetivo:} construir a API de uma biblioteca pública — autores, livros
e empréstimos — com Prisma, migrações e consultas tipadas.

\textbf{Critérios de aceite:}
\begin{itemize}[leftmargin=*, itemsep=0.2em]
  \item Schema com 3 tabelas relacionadas e pelo menos 2 migrações no histórico;
  \item CRUD de autores e livros, e registro de empréstimos (livro $\leftrightarrow$ leitor, com \texttt{dataDevolucao});
  \item Endpoint de busca de livros por título/autor (contém, case-insensitive);
  \item Endpoint ``livros emprestados'' sem N+1 (verifique com logs de query);
  \item Seed\index{seed} com 10 autores e 30 livros reais (clássicos da computação e literatura);
  \item Exercício bônus: refazer o mesmo projeto com Drizzle e comparar as duas experiências num README.
\end{itemize}
\end{caixaprojeto}

\begin{caixadica}
Migrações em produção exigem cuidado: rode \texttt{prisma migrate deploy}
no deploy (capítulo~\ref{cap:cicd}), nunca em dev. E antes de migrar em
produção, \emph{backup} — o capítulo~\ref{cap:postgresql} te ensinou
\texttt{pg\_dump}.
\end{caixadica}

\section{Exercícios de fixação}

\begin{exercicios}
  \item O que é uma migração e por que versioná-la?
  \item Modele no Prisma: \texttt{Usuario} 1--N \texttt{Post} 1--N \texttt{Comentario}.
  \item Escreva uma query Prisma que retorna posts com autor e comentários de um usuário específico.
  \item Escreva a mesma query em Drizzle usando \texttt{innerJoin}.
  \item Explique o problema N+1 com um exemplo concreto e as duas soluções.
  \item O que faz \texttt{prisma migrate deploy} e quando usá-lo?
\end{exercicios}

\section{Desafios}

\begin{desafios}
  \item \textbf{Migração reversível}: crie uma migração, reverta-a (\texttt{migrate dev --rollback} ou equivalente) e documente o que acontece com os dados.
  \item \textbf{Transação com Prisma}: implemente a transação ``registrar empréstimo + marcar livro como emprestado'' com \texttt{prisma.\$transaction}.
  \item \textbf{Busca com paginação}: adicione \texttt{?pagina=\&limite=} na busca de livros e retorne o total.
  \item \textbf{Comparativo Drizzle}: refaça o CRUD com Drizzle e escreva o README comparando DX, tipos e queries.
\end{desafios}

\section{Exercícios de nível sênior}

\begin{seniores}
  \item \textbf{Índices via schema}: adicione índices no Prisma/Drizzle para as colunas de busca e comprove a melhoria com \texttt{EXPLAIN ANALYZE}.
  \item \textbf{Data de devolução}: implemente a regra ``livro atrasado'' com comparação de datas no banco (não no Node) e um endpoint de relatório de atrasos.
  \item \textbf{Soft delete}: implemente exclusão lógica (\texttt{deletedAt}) e explique os trade-offs vs.\ exclusão física.
  \item \textbf{Seed idempotente}: escreva um seed que roda N vezes sem duplicar dados (\texttt{upsert}).
\end{seniores}

\section{Armadilhas comuns}

\begin{itemize}[leftmargin=*, itemsep=0.4em]
  \item Fazer \texttt{include} em cascata sem necessidade (traga o mínimo);
  \item N+1 queries silenciosas (ligue o logging de queries em dev);
  \item Alterar o banco fora das migrações (schema e código dessincronizam);
  \item Colocar \texttt{DATABASE\_URL} no código (use \texttt{.env}, capítulo~\ref{cap:seguranca});
  \item Migrar em produção sem backup/teste (capítulo~\ref{cap:cicd});
  \item \texttt{prisma generate} esquecido após mudar o schema (client desatualizado).
\end{itemize}

\section{Resumo do capítulo}

\begin{itemize}[leftmargin=*, itemsep=0.3em]
  \item ORMs traduzem tabelas em tipos e geram queries seguras;
  \item Migrações versionam o schema como o Git versiona o código;
  \item Prisma: abstração e DX; Drizzle: SQL-first e leve — saiba escolher;
  \item N+1 é o anti-padrão nº 1; \texttt{include}/\texttt{JOIN} resolvem;
  \item Log de queries ligado em dev revela problemas silenciosos;
  \item Projeto entregue: API de biblioteca com migrações e seed.
\end{itemize}

\section{Referências oficiais para aprofundar}

\begin{itemize}[leftmargin=*, itemsep=0.3em]
  \item Documentação do Prisma: \url{https://www.prisma.io/docs}
  \item Documentação do Drizzle: \url{https://orm.drizzle.team}
  \item Prisma vs Drizzle (comparação oficial): \url{https://www.prisma.io/docs/orm/more/comparisons/prisma-and-drizzle}
\end{itemize}
